\documentclass[a4paper]{article}
\usepackage{amsfonts}
\usepackage{amsmath}
\usepackage{amssymb}
\usepackage{graphicx}   

\begin{document}
  
\noindent \large Auteur: Etienne El Hachem \\
\noindent \large Studierichting: Informatica\\
\noindent \large Oefeningenreeks 2D nr. 10.7\\

\medskip

\normalsize

$\lim\limits_{x\rightarrow -2\sqrt{5}} \dfrac{\sqrt{5}x^2+9x-2\sqrt{5}}{2x^2+3x\sqrt{5}-10} = \dfrac{0}{0}$ \\

$\sqrt{5}x^2+9x-2\sqrt{5}$: \\

\begin{tabular}{c|cc|c}
& $\sqrt{5}$ & $9$ & $-2\sqrt{5}$ \\
$-2\sqrt{5}$ & & $-10$ & $2\sqrt{5}$ \\ \hline
& $\sqrt{5}$ & $-1$ & $0$
\end{tabular}
$\Rightarrow (x+2\sqrt{5})(\sqrt{5}x - 1)$ \\


$2x^2+3x\sqrt{5}-10$: \\

\begin{tabular}{c|cc|c}
& $2$ & $3\sqrt{5}$ & $-10$ \\
$-2\sqrt{5}$ & & $-4\sqrt{5}$ & $10$ \\ \hline
& $2$ & $-\sqrt{5}$ & $0$
\end{tabular}
$\Rightarrow (x + 2\sqrt{5})(2x - \sqrt{5})$ \\


$\Rightarrow \lim\limits_{x\rightarrow -2\sqrt{5}} \dfrac{(x+2\sqrt{5})(\sqrt{5}x - 1)}{(x + 2\sqrt{5})(2x - \sqrt{5})}$ \\

$= \lim\limits_{x\rightarrow -2\sqrt{5}} \dfrac{(\sqrt{5}x - 1)}{(2x - \sqrt{5})}$ \\

$= \dfrac{11}{5\sqrt{5}}$

\begin{thebibliography}{999}
%\bibitem{a} Referentie
\end{thebibliography}
\end{document} 
